\documentclass[11pt]{article}
\usepackage[margin=1in]{geometry}
\usepackage{amsmath,amssymb,amsthm,hyperref}
\newtheorem{theorem}{Theorem}
\newtheorem{corollary}[theorem]{Corollary}
\newcommand{\one}{\boldsymbol 1}
\newcommand{\R}{\mathbb R}
\title{A flat, optimal balanced cut can flip after one spectrally vanishing edge doubling}
\author{Partitioning limits research session}
\date{15 November 2026}
\begin{document}
\maketitle

This note gives an explicit obstruction, not a claim of historical novelty.
The Fiedler vectors refer to the combinatorial Laplacian and unit vertex
masses. All graphs below are dense, simple, positively weighted graphs.
No bounded-degree or unweighted conclusion is asserted.

\begin{theorem}\label{thm:main}
For each integer $n\ge4$ divisible by four there are two connected
positively weighted simple graphs $G_-,G_+$ on the same $n$ vertices such that:
\begin{enumerate}
\item $G_+$ is obtained from $G_-$ by doubling the weight of exactly one
existing edge, from $2/3$ to $4/3$.
\item Both Fiedler eigenvalues are simple. Their unit Fiedler vectors
$v_-,v_+$ are exactly orthogonal, have no zero entries, and satisfy
\[
 \frac1{\sqrt{5n}}\le |v_\pm(i)|\le\frac3{\sqrt n}
 \qquad(1\le i\le n).
\]
\item Their sign cuts are the respective unique minimum bisections, up
to complementation, and their distance after identifying complements is
exactly $n/2$ vertices.
\item If $A=L(G_-)$ and $B=L(G_+)$, then
\[
 0\preceq B-A=\frac23bb^T,\qquad
 \|A^{\dagger/2}(B-A)A^{\dagger/2}\|
 \le\frac4{3n-5}\longrightarrow0.
\]
For every fixed $0<\alpha\le1$, the relative fractional correction satisfies
\[
 0\preceq A^{\dagger\alpha/2}(B^\alpha-A^\alpha)
                    A^{\dagger\alpha/2}
 \preceq
 \left[\left(1+\frac4{3n-5}\right)^\alpha-1\right]P,
\]
where $P=I-\one\one^T/n$ is the projection off constants.
\item Every bisection is a $(1+2/n)$-approximate minimum bisection
in both $G_-$ and $G_+$. The linear span of all these normalized
centered cut vectors is the entire $(n-1)$-dimensional space $\one^\perp$.
More quantitatively, a subspace $W$ whose squared-distance error is at
most $\eta$ for every one of these unit cut vectors must have dimension
at least $(1-\eta)(n-1)$.
\end{enumerate}
\end{theorem}

\begin{proof}
Split the vertices into four equal nonempty blocks and define unit vectors
constant on these blocks by
\[
 f=n^{-1/2}(+,-,-,+),\qquad
 h=n^{-1/2}(+,+,-,-).
\]
Here each sign is repeated $n/4$ times. The vectors $\one/\sqrt n,f,h$
are orthonormal. Put
\[
 L_0=nI-\one\one^T-ff^T-hh^T=nP-ff^T-hh^T.
\]
Its off-diagonal entries are $-w_{ij}$ with
\[
 w_{ij}=1+f_if_j+h_ih_j\in\{1-2/n,1,1+2/n\}.
\]
They are strictly positive, and $L_0\one=0$, so $L_0$ is the
Laplacian of a positively weighted complete graph. Choose $u$ in block
one and $v$ in block two. For $b=e_u-e_v$ we have
\[
 w_{uv}=1,\quad b^Th=0,\quad a:=b^Tf=2/\sqrt n,
 \quad r:=b-af\perp\{\one,f,h\}.
\]
In particular $\|r\|^2=2-a^2>0$. Define
\[
 A=L_0-\tfrac13bb^T,\qquad B=L_0+\tfrac13bb^T.
\]
These alter only $w_{uv}$, to $2/3$ and $4/3$ respectively.

On $\one^\perp$, $L_0$ has eigenvalue $n-1$ on
$\operatorname{span}\{f,h\}$ and eigenvalue $n$ on its orthogonal
complement. Let $\delta>0$ be the unique solution
\[
 3\delta^2+\delta-a^2=0,
 \qquad
 \delta=\frac{\sqrt{1+48/n}-1}{6}\le a^2=4/n.
\]
Direct multiplication gives
\[
 A(af+\delta b)=(n-1-\delta)(af+\delta b).
\]
For example, this follows by using $L_0f=(n-1)f$,
$L_0b=nb-af$, and $b^T(af+\delta b)=a^2+2\delta$.
The coefficient of $b$ is correct precisely because
$3\delta^2+\delta-a^2=0$.
Since a negative rank-one perturbation can put at most one eigenvalue
strictly below the previous minimum $n-1$, this is the unique Fiedler
eigenvalue of $A$. Alternatively, this follows from the explicit
two-dimensional restriction to $\operatorname{span}\{f,r\}$.

The vector $x=af+\delta b$ has exactly the signs of $f$: away from
$u,v$ it is a positive multiple of $f$, and at $u,v$ the added term
reinforces their respective signs. Also
\[
 \|x/a\|^2=\|f+(\delta/a)b\|^2
 =1+2\delta+2\delta^2/a^2
 \le1+4\delta\le5.
\]
Every coordinate of $x/a$ has magnitude at least $1/\sqrt n$.
At either endpoint its magnitude is
\[
 \frac1{\sqrt n}+\frac\delta a
 =\frac{1+n\delta/2}{\sqrt n}\le\frac3{\sqrt n};
\]
all other coordinates have magnitude exactly $1/\sqrt n$.
Normalizing $x$ proves the asserted bounds for $v_-$.

The vector $h$ is an eigenvector of $B$ with eigenvalue $n-1$.
For every $y\perp\one$,
\[
 y^T(B-(n-1)I)y
 =\|\Pi_{\{\one,f,h\}^\perp}y\|^2+\tfrac13(b^Ty)^2.
\]
Equality requires $y\in\operatorname{span}\{f,h\}$ and $b^Ty=0$,
which force $y\in\operatorname{span}\{h\}$. Thus $h$ is the unique
Fiedler direction of $B$. It is exactly orthogonal to $x$ because
both $f$ and $b$ are orthogonal to $h$. Set $v_+=h$.

Write a bisection as a balanced sign vector
$s\in\{\pm1\}^n$, $s^T\one=0$. Its cut weight for $L_0$ is
\[
 \frac14s^TL_0s
 =\frac{n^2}{4}-\frac{(f^Ts)^2+(h^Ts)^2}{4}
 \ge q:=\frac{n(n-1)}4.
\]
Bessel's inequality gives the lower bound. Equality forces
$s\in\operatorname{span}\{f,h\}$. Writing
$s=\sqrt n(cf+dh)$, its values on the four blocks are
$c+d,-c+d,-c-d,c-d$. All have magnitude one only when $cd=0$,
$c^2+d^2=1$. Hence equality occurs only at
$s=\pm\sqrt n f$ or $s=\pm\sqrt n h$.
The edge $uv$ crosses the $f$-cut and does not cross the $h$-cut.
Consequently the minimum cut for $A$ is uniquely the $f$-cut with
weight $q-1/3$, and the minimum for $B$ is uniquely the $h$-cut
with weight $q$. These cuts differ on precisely blocks two and four,
so their complement-invariant distance is $n/2$.

Since $L_0\succeq(n-1)P$ and $\|bb^T\|=2$,
\[
 A\succeq(n-5/3)P.
\]
Therefore
\[
 \tau:=\left\|A^{\dagger/2}\frac23bb^TA^{\dagger/2}\right\|
 =\frac23b^TA^\dagger b\le\frac{4}{3n-5}.
\]
On $\one^\perp$, $A\preceq B\preceq(1+\tau)A$.
The operator monotonicity of $x^\alpha$ for $0<\alpha\le1$
gives the stated relative fractional bound by congruence. All powers
vanish on constants.

Every bisection has $L_0$ cut weight at most $n^2/4$.
Its $A$ cut weight lies between $q-1/3$ and $n^2/4$, while its
$B$ cut weight lies between $q$ and $n^2/4+1/3$.
For $n\ge4$,
\[
 \frac{n/4+1/3}{q-1/3}\le\frac2n,
\]
because this inequality is equivalent to $(3n+2)(n-4)\ge0$.
The same bound follows for $B$, whose denominator $q$ is larger.
Finally balanced sign vectors span $\one^\perp$: the difference of
two bisections that exchange vertices $i,j$ is $2(e_i-e_j)$,
and these differences span $\one^\perp$.
For the quantitative assertion, choose a balanced sign vector $s$
uniformly at random and set $z=s/\sqrt n$. Symmetry and $s^T\one=0$
give $\mathbb E zz^T=P/(n-1)$. Thus for every subspace $W$,
\[
 \mathbb E\|\Pi_Wz\|^2=\frac{\operatorname{tr}(\Pi_WP)}{n-1}
 \le\frac{\dim W}{n-1}.
\]
If all squared errors are at most $\eta$, this expectation is at least
$1-\eta$, proving the claim.
\end{proof}

\begin{corollary}
There is no modulus tending to zero that bounds either the angle
between unique Fiedler directions or the complement-invariant fraction
of vertices changed by their optimal balanced sign cuts in terms only
of the relative fractional perturbation norm of one existing-edge
doubling. This remains true under uniform lower and upper bounds
of order $n^{-1/2}$ on all normalized Fiedler coordinates, and with
algebraic connectivity tending to infinity.
\end{corollary}

The exact ordinary relative perturbation norm is also available.
Sherman--Morrison on $\one^\perp$ gives
\[
 r_0=b^TL_0^\dagger b
 =\frac{a^2}{n-1}+\frac{2-a^2}{n}
 =\frac2n+\frac4{n^2(n-1)},\qquad
 \tau=\frac{(2/3)r_0}{1-r_0/3}.
\]
Thus $n\tau\to4/3$. The obstruction has vanishing eigengap between
the first two nontrivial eigenvalues; it does not contradict classical
gap-dependent perturbation theorems. Its force is that localization,
rounding near zero, disconnectedness, and nonoptimal Fiedler cuts are
all absent from the example.

The last assertion about all nearoptimal cuts is elementary in dense
graphs: almost all balanced cuts have approximately the same dominant
$n^2/4$ cost. It is included to prevent a spurious low-dimensional
all-nearoptimal-cut claim from being inferred from a rank-one update.
Neither this corollary nor the basic eigenspace-tail estimate should be
advertised as a new structural theorem without a dedicated priority check.
\end{document}
